\section{REFERENCIAL TEÓRICO PRELIMINAR}

\subsection{A Maternidade Solo no Brasil e a Economia do Cuidado}
A maternidade aplica um peso desigual a mulheres e homens na sociedade, desde a gestação até a criação e suporte financeiro à criança. Conforme afirmado por \textcite{sousa2025}, os homens têm a designação prioritária à esfera produtiva, enquanto a esfera reprodutiva recai sobre as mulheres. Ao compararmos essa afirmação com o dado apresentado também por \textcite[28]{sousa2025}, de que quase metade das mulheres no contexto analisado (47\%) não recebe pensão alimentícia ou ajuda financeira dos genitores, e também com o cenário descrito por \textcite[5]{souza2024}, onde a mulher está sujeita a uma “tripla penalidade” (tempo escasso, renda insuficiente e ausência de redes de apoio efetivas), o panorama para as responsabilidades de criação dos filhos se torna ainda menos otimista.

\subsection{O Paradigma da Tecnologia Social e Redes Colaborativas (RTS)}
Uma das ferramentas mais poderosas à disposição para amenizar, em partes, os problemas enfrentados por mães solo são as interações sociais, ainda que virtuais. Como observado por \textcite[694]{brito2022}:
\begin{quote}
As mulheres utilizam suas próprias experiências sobre a maternidade na intenção de ajudar outras mulheres, de modo a criar uma rede de apoio (social e emocional), especialmente para as primíparas [mulheres na primeira gestação], proporcionando informação, incentivo e empatia.
\end{quote}
O contato direto com outros indivíduos que passaram ou passam por situações análogas ajuda a criar um sentimento de comunidade e pertencimento. Além disso, cria um espaço seguro para desabafo e ressignificação de experiências angustiantes, como descrito por \textcite[6]{araujo2025}.

Para além disso, a criação de uma plataforma centralizada que facilita doações tem impacto direto e imediato na qualidade de vida das mães solo inseridas nos demográficos de vulnerabilidade financeira, pois supre, ao menos parcialmente, a lacuna deixada pela falta de renda apropriada e de apoio financeiro por parte dos genitores. Como observam \textcite[9]{araujo2025}, a renda vinda de programas como o Bolsa Família, “ainda que seja essencial para as despesas básicas do lar, não é suficiente”, o que torna o apoio de terceiros decisivo para que essas mães tenham acesso a produtos infantis como enxoval, fraldas, fórmulas ou móveis que elas não conseguiriam comprar por conta própria.

\subsection{Arquitetura Hexagonal (Ports \& Adapters) e Blindagem de Domínio}
Podemos observar, na forma como os projetos estudantis são comumente estruturados, certa tendência a pontos de atrito arquitetônico que geram, a depender do tamanho e da complexidade do projeto, dificuldade de escalabilidade e manutenibilidade. Segundo \textcite{deluca2025}, escolhas inadequadas na estrutura do código geram dívida técnica e queda no desempenho e na confiabilidade, e a partir disso podemos entender a relevância que toma o cuidado e a atenção com esses detalhes de projeto.

Ao nos debruçarmos na tarefa de entender onde esses erros mais aparecem, notamos que há uma predominância na questão tangente ao acoplamento, ou seja, não existe uma divisão clara de qual a responsabilidade de cada componente. Isso acaba levando a cenários como o descrito por \textcite[9]{deluca2025}, em que 25\% dos projetos de alunos analisados colocavam chamadas de banco de dados (DAO) diretamente dentro da interface gráfica (GUI). Este trabalho aborda esse cenário com a aplicação formal e, até certo ponto, rígida dos princípios S.O.L.I.D. e dos padrões de projeto chamados de Arquitetura Hexagonal \cite{rayarao2026}, para diminuir o acoplamento e facilitar tanto a leitura e manutenção do código quanto futuras expansões do escopo geral da aplicação.

\subsection{Privacidade por Padrão (Privacy by Design) e a LGPD}
Observemos a questão da privacidade por dois prismas que podem parecer distintos, mas chegam ao mesmo ponto. O primeiro refere-se ao cuidado necessário com os dados de pessoas em situações de vulnerabilidade econômica e social. Como visto em \textcite[11]{araujo2025}, são frequentes entre mães solo os relatos de:
\begin{quote}
violência doméstica dos tipos física, psicológica, sexual, patrimonial e moral perpetradas pelos ex-companheiros
\end{quote}
um dado que pede maior atenção e delicadeza, visto que, ao expor dados sensíveis, corre-se o risco de entregar informações a terceiros mal-intencionados, inclusive ex-parceiros que as usuárias podem estar evitando.

O outro ponto de vista a tratar de segurança é a Lei Geral de Proteção de Dados (LGPD), que legisla sobre a forma como os dados são armazenados e expostos, bem como sobre a clareza para o usuário quanto ao motivo da coleta. Essas diretrizes podem ser, de forma geral, entendidas em \textcite[4]{moura2024}, quando apontam que ``a coleta de dados deve estar associada a uma necessidade ou finalidade específica, bem definida, deve ser clara e estar em conformidade com as diretrizes vigentes na legislação''.